We run a small, CPU-only comparison of reimplemented operator-learning models on the 1D viscous Burgers equation ($\nu=0.01$, periodic): a Fourier neural operator (FNO), a DeepONet, an MLP and a CNN, each trained on 400 pairs of initial condition and solution at $t=1$ on a 128-point grid, with five seeds. The FNO reaches a test relative $L_2$ error of about $0.005$, versus $0.083$ for the MLP and about $0.22$ for DeepONet and the CNN, under an identical 100-epoch budget; DeepONet and the MLP remain far behind the FNO with 5--20$\times$ more epochs (three seeds). Our plain DeepONet thus does not reproduce the near-parity with the FNO reported in the literature for simple settings, and its error is a lower bound on DeepONet quality. Evaluated zero-shot on 256- and 512-point grids, the FNO and the CNN accept the finer input natively; the FNO error is unchanged while the CNN error grows by factors of 2.55 and 3.27. MLP and DeepONet are evaluated by stride-subsampling the input, so their unchanged error is a property of this rule and not evidence of resolution invariance. Because the initial conditions are band-limited and all grids resolve the solution, this test is a clean but easy one. On three seeds the FNO error shows no measurable sensitivity to the number of Fourier modes from 8 to 32 or to removing the grid channel, and falls steadily as the training set grows from 100 to 400 pairs, with the gap to DeepONet persisting at every size (at a fixed number of epochs, so fewer pairs also mean fewer gradient steps). The study is small: one PDE, one data distribution, tuned learning rates only, and baselines that may be under-trained.
